\documentclass[12pt,a4paper]{article}

% Math
\usepackage{amsmath,amssymb,amsfonts,amsthm,mathtools}
\usepackage{physics}
\usepackage{bm}

% Tables
\usepackage{array,booktabs,multirow,multicol,longtable,tabularx}
\usepackage{colortbl,xcolor}

% Figures
\usepackage{graphicx,float,subcaption}

% TikZ
\usepackage{tikz,pgfplots}
\pgfplotsset{compat=1.18}

% Layout
\usepackage[a4paper,margin=1.5cm]{geometry}
%\usepackage{anysize}
%\marginsize{1.5cm}{1.5cm}{1.5cm}{1.5cm}

% Utilities
\usepackage{enumitem}
\usepackage{makecell}
\usepackage{hyperref}

% Vietnamese
\usepackage[utf8]{inputenc}
\usepackage[T5]{fontenc}
\usepackage[vietnamese]{babel}

\title{\textbf{Introduction to \LaTeX}}
\author{VyNHT}
\date{\today}

\begin{document}
	
	\maketitle 
	% \begin{center}
		% 	\LARGE{\textbf{Introduction to \LaTeX}}\\ \vspace{2mm}
		% 	\large{VyNHT}
		% \end{center}
	
	The Fourier series of a periodic function \(f(x)\) can be represented as
	\[f(x) = a_0 + \sum_{n=1}^{\infty} \left( a_n \cos\left(\frac{2\pi n x}{T}\right) + b_n \sin\left(\frac{2\pi n x}{T}\right) \right)\]
	where \(a_n,\ b_n\) are the Fourier coefficients, calculated as follows
	\[a_n = \frac{2}{T} \displaystyle\int_0^T f(x) \cos\left(\frac{2\pi n x}{T}\right)\; \mathrm{d}x,\quad b_n = \frac{2}{T} \displaystyle\int_0^T f(x) \sin\left(\frac{2\pi n x}{T}\right)\; \mathrm{d}x.\]

\end{document}